\chapter{Conclusão}
\label{cap:conclusao}

Este relatório documentou a primeira etapa do Projeto de Formatura, dedicada ao desenvolvimento 
e à avaliação de uma estratégia de controle distribuído em dupla camada para BESS, voltada à 
regulação de tensão em redes de distribuição de média tensão com alta penetração de geração 
fotovoltaica. A seguir, retomam-se os objetivos específicos definidos na 
Seção~\ref{sec:objetivos}, confrontando-os com os resultados obtidos.

A revisão da literatura (Capítulo~\ref{cap:fundamentacao}) sistematizou as arquiteturas de 
controle e as principais formulações distribuídas, identificando os algoritmos de consenso 
como a base predominante da camada de coordenação e fundamentando as escolhas metodológicas 
do trabalho. O ambiente de cossimulação Python--OpenDSS (Capítulo~\ref{cap:ambiente}) foi 
desenvolvido e validado na rede IEEE~33 barras, estabelecendo a infraestrutura para a execução 
de fluxos de potência iterativos em séries temporais. A geração de cenários por Monte Carlo 
(Capítulo~\ref{cap:cenarios}) produziu um conjunto de 5500~cenários sobre a rede IEEE~34 
barras, contemplando onze níveis de penetração e a incerteza de irradiância, alocação de 
recursos e demanda. As duas camadas de controle (Capítulo~\ref{cap:controle}) foram 
implementadas e integradas, e seu desempenho foi avaliado de forma pareada 
(Capítulo~\ref{cap:resultados}).

Quanto à avaliação de desempenho, três conclusões se destacam. Primeiro, nas condições 
simuladas, o desafio dominante de regulação de tensão é a sobretensão diurna, que cresce 
sistematicamente com a penetração fotovoltaica, enquanto a subtensão é satisfatoriamente 
mitigada pelos dispositivos convencionais da rede, à exceção da barra de extremidade do 
ramal de 4,16~kV. Segundo, o controle primário por curva volt-var mostrou-se eficaz: 
reduziu de forma expressiva o número médio de barras em sobretensão (de cerca de dezessete 
para quatro em 100\% de penetração) e comprimiu a dispersão entre cenários, tornando o 
desempenho da rede mais previsível, sem causar corte de geração. Terceiro, o controle 
secundário por consenso, na rede e configuração estudadas, não reduziu as violações 
remanescentes ao primário, apresentando inclusive efeito adverso pontual; tal resultado 
decorre da concentração espacial dos BESS no alimentador, que limita o alcance da 
redistribuição de potência sobre as barras críticas próximas à subestação.

O resultado negativo da camada de consenso, longe de invalidar a abordagem, delimita suas 
condições de eficácia e orienta os próximos passos. Cabe ressaltar dois ajustes em relação ao 
plano de trabalho, ambos justificados ao longo do texto: a adoção da biblioteca 
OpenDSSDirect.py em lugar da interface COM, e a substituição da lógica \emph{fuzzy} pela curva 
volt-var normatizada, estendida também aos inversores FV.

\section{Trabalhos futuros}
\label{sec:futuro}

Os resultados obtidos delineiam os próximos passos da pesquisa, a serem desenvolvidos na 
segunda etapa do Projeto de Formatura:

\begin{itemize}
  \item aprimorar os modelos de FV e BESS, em particular incorporando o rastreamento do estado de carga (SoC) das baterias, hoje modeladas como elementos de carga equivalente;
  \item reavaliar o desempenho da rede na ausência dos reguladores automáticos de tensão, isolando a contribuição do controle distribuído;
  \item revisar a alocação e o dimensionamento do armazenamento, incluindo mais unidades ou uma rede mais favorável, de modo a criar condições adequadas à atuação da camada de consenso, e definir novos critérios de desempenho, como redução de perdas, suporte de reativos e mitigação de desequilíbrio;
  \item estender a solução a redes trifásicas desequilibradas, analisando o impacto da assimetria de fases sobre o algoritmo de controle;
  \item incorporar modelos de degradação e ciclo de vida das baterias ao laço de controle; e
  \item aplicar a metodologia a uma rede de distribuição real, aproximando a avaliação das condições operacionais efetivas.
\end{itemize}
